\documentclass[10pt, twocolumn, letterpaper]{article}

\usepackage[utf8]{inputenc}
\usepackage{graphicx}
\usepackage{amsmath}
\usepackage{amssymb}
\usepackage{booktabs}
\usepackage{hyperref}
\usepackage{geometry}
\usepackage{abstract}

\geometry{left=1.5cm, right=1.5cm, top=2cm, bottom=2cm}

\title{\textbf{SENSERA: A Multi-Stage AI Framework for Volumetric Lung Nodule Detection and Radiogenomic Biomarker Prediction}}

\author{
    Roberto Hordoan, Razvan Bocra, Marius Boroica, Mihai Deaconu \\
    \textit{Department of Computer Science} \\
    \textit{Projects In Silico Diagnostics}
}

\date{\today}

\begin{document}

\maketitle

\begin{abstract}
Lung cancer remains the leading cause of cancer-related mortality worldwide, necessitating advanced tools for early detection and non-invasive characterization. This study presents SENSERA, a comprehensive AI framework integrating a state-of-the-art 3D segmentation pipeline with a radiogenomic biomarker prediction model. Our segmentation approach leverages VISTA-3D, a foundation model trained on 11,000 CT scans, augmented by a "Shadow Mining" 3D ResNet-18 classifier for false positive reduction. This cascade achieves a volumetric tumor Dice score of 0.7342 and drastically reduces false positives to 0.58 per scan. Parallelly, our radiogenomics module employs a late fusion strategy combining deep learning features (DenseNet121) with handcrafted radiomics (XGBoost) to predict critical biomarkers from CT imaging. We report Area Under the Curve (AUC) scores of 0.87 for EGFR and 0.86 for KRAS mutations on public datasets, demonstrating the system's potential as a non-invasive "virtual biopsy" tool.
\end{abstract}

\section{Introduction}
The accurate detection of pulmonary nodules and the subsequent characterization of their molecular properties are critical steps in lung cancer management. While Low-Dose Computed Tomography (LDCT) is the standard for screening, it suffers from high false-positive rates and cannot inherently resolve molecular subtypes (e.g., EGFR, KRAS mutations) without invasive biopsies. SENSERA addresses these dual challenges through two interconnected pipelines: (1) a hierarchical segmentation system that effectively detects relevant nodules while suppressing artifacts, and (2) a multi-modal radiogenomic framework that infers genetic information directly from imaging data.

\section{Methodology}

\subsection{3D Segmentation and Detection Pipeline}
Our detection framework adopts a coarse-to-fine strategy, utilizing a powerful foundation model for initial candidate generation followed by a specialized classifier for refinement.

\subsubsection{Stage 1: Foundation Model (VISTA-3D)}
The primary segmentation engine is built upon \textbf{VISTA-3D} (NVIDIA/MONAI), a large-scale 3D Transformer architecture comparable to the Segment Anything Model (SAM).
\begin{itemize}
    \item \textbf{Pretraining:} The model was pretrained on a massive corpus of approximately 11,000 CT scans spanning 127 anatomical classes, enabling robust generalization across diverse patient morphologies.
    \item \textbf{Inference:} We utilize the model's auto-segmentation capabilities to generate initial nodule candidates, capturing high-sensitivity regions of interest (ROIs).
\end{itemize}

\subsubsection{Stage 2: "Shadow Mining" False Positive Reduction}
To address the high sensitivity but potentially low precision of the foundation model, we implement a secondary "Shadow Mining" classifier.
\begin{itemize}
    \item \textbf{Architecture:} A modified 3D ResNet-18 serves as the discriminator.
    \item \textbf{Input Configuration:} The model processes a 2-channel volumetric input: the raw CT attenuation data (Hounsfield Units) and the probability map generated by Stage 1.
    \item \textbf{Shadow Mining Mechanism:} We employ an ensemble-like mining strategy where five separate VISTA models are trained on distinct cross-validation splits. These models are used to mine "hard" false positives (FPs) from unseen validation folds, creating a challenging training set for the classifier.
    \item \textbf{Optimization:} The classifier is trained using \textbf{Focal Loss} to address class imbalance and \textbf{Mixup} data augmentation to regularize the decision boundary and prevent overfitting on texture artifacts.
\end{itemize}

\subsection{Radiogenomics and Biomarker Prediction}
Our biomarker prediction module predicts genetic mutations (EGFR, KRAS) and histological subtypes (TTF-1, CK7) by synthesizing deep visual features with quantitative radiomics.

\subsubsection{Data Strategy and Pretraining}
We leverage a synergy of public datasets (TCGA-LUAD, NSCLC Radiogenomics) and a private cohort of 92 patients.
\begin{itemize}
    \item \textbf{Bio-Specific Pretraining:} A 3D DenseNet121 is initialized by pretraining on RNA-seq protein expression data (specifically MK167 proliferation markers) to learn latent representations of tumor aggressiveness.
    \item \textbf{Fine-Tuning:} The pretrained backbone is fine-tuned on the private cohort using Immunohistochemistry (IHC) stain records as ground truth.
\end{itemize}

\subsubsection{Late Fusion Architecture}
We employ a late fusion ensemble to combine complementary feature streams:
\begin{itemize}
    \item \textbf{Deep Branch:} The fine-tuned DenseNet121/ResNet extracts high-level semantic features representing spatial morphology and intratumoral heterogeneity.
    \item \textbf{Radiomics Branch:} An XGBoost classifier is trained on handcrafted features (texture, shape, intensity statistics) extracted via PyRadiomics.
    \item \textbf{Fusion:} The probability outputs of both branches are fused using a Logistic Regression meta-learner to produce the final diagnostic prediction.
\end{itemize}

\section{Experimental Results}

\subsection{Segmentation Performance}
The integration of the "Shadow Mining" ResNet classifier yielded State-of-the-Art (SOTA) improvements over the standalone foundation model. As shown in Table~\ref{tab:seg_results}, precision was more than doubled, and the false positive rate was reduced by nearly 80\%.

\begin{table}[h]
\centering
\caption{Segmentation Performance Comparison}
\label{tab:seg_results}
\resizebox{\columnwidth}{!}{%
\begin{tabular}{lcc}
\toprule
\textbf{Metric} & \textbf{VISTA Only} & \textbf{VISTA + ResNet (Ours)} \\
\midrule
Tumor Dice (Volumetric) & 0.6260 & \textbf{0.7342} \\
Recall (Sensitivity) & \textbf{0.9447} & 0.9213 \\
Precision & 0.2745 & \textbf{0.6483} \\
FPs per Scan & 2.76 & \textbf{0.58} \\
Lesion-wise Dice (TPs) & 0.7858 & \textbf{0.8102} \\
\bottomrule
\end{tabular}%
}
\end{table}

While Recall dropped slightly (approx. 2.3\%), the dramatic increase in Precision (from 0.27 to 0.64) underscores the effectiveness of the Shadow Mining approach for clinical utility, where false alarms are a major bottleneck.

\subsection{Biomarker Prediction}
The Late Fusion strategy consistently outperformed individual modalities (Deep Learning only or Radiomics only) across all tested mutations.

\subsubsection{Public Genomic Datasets}
Results on the combined TCGA-LUAD and NSCLC Radiogenomics test sets are presented below.

\begin{table}[h]
\centering
\caption{Biomarker Prediction (Public Data)}
\label{tab:biomarker_public}
\resizebox{\columnwidth}{!}{%
\begin{tabular}{lccc}
\toprule
\textbf{Target} & \textbf{DL Only} & \textbf{Radiomics (XGB)} & \textbf{Late Fusion (Best)} \\
\midrule
EGFR Mutation & 0.73 AUC & 0.82 AUC & \textbf{0.87 AUC} \\
KRAS Mutation & 0.71 AUC & 0.81 AUC & \textbf{0.86 AUC} \\
\bottomrule
\end{tabular}%
}
\end{table}

\subsubsection{Private Dataset Validation}
On our internal cohort of 92 patients with IHC validation, the model demonstrated strong generalization capabilities:
\begin{itemize}
    \item \textbf{TTF-1 Prediction:} 0.80 AUC (Late Fusion)
    \item \textbf{CK7 Prediction:} 0.76 AUC (Late Fusion)
\end{itemize}

\section{Discussion and Next Steps}
The SENSERA pipeline demonstrates that combining foundation models with specialized false-positive reduction networks can achieve clinical-grade detection performance. Furthermore, our results validate the "Virtual Biopsy" hypothesis, showing that non-invasive imaging can predict specific mutations with high accuracy (AUC $> 0.85$).

\subsection{Future Work}
\begin{enumerate}
    \item \textbf{Data Preprocessing Refinement:} We aim to detail and integrate advanced filtering techniques, such as the LUNA16 filter, to further clean input data and standardize voxel spacing across heterogeneous datasets.
    \item \textbf{Clinical Integration:} Future phases will focus on deploying the "Shadow Mining" loop in an active learning setup, allowing radiologists to iteratively improve the FP reduction model.
    \item \textbf{Multicenter Validation:} Expanding the private cohort validation to external institutions to confirm the robustness of the TTF-1 and CK7 predictors.
\end{enumerate}

\section{Conclusion}
SENSERA establishes a robust baseline for AI-driven lung cancer diagnostics, bridging the gap between anatomical detection and molecular characterization. The high precision in detection and strong correlation with genomic markers suggest a viable path toward reduced biopsy rates and personalized treatment planning.

\end{document}
